8.1 Purpose of the Simulation Environment
The BED v3 simulation environment serves as a computational testbed for validating incentive policies before committing to costly field deployment. As stated in your uploaded document:

“The BED v3 simulation environment serves as a computational testbed for validating incentive policies…”

The simulation allows researchers to:

evaluate policy performance under controlled conditions

test habit dynamics and behavioral response

quantify reward spend and cost‑effectiveness

perform sensitivity analysis

conduct power analysis for field experiments

ensure fairness and budget constraints are respected

This chapter defines the full simulation architecture, behavioral model, experimental conditions, and evaluation metrics.

8.2 Simulation Overview
The simulation is a discrete‑time, agent‑based environment that mirrors the real BED v3 prototype:

Each agent represents a consumer.

Each day represents a decision opportunity.

Incentives influence event probability.

Events update habit state.

Policies determine reward multipliers.

Metrics track behavioral and economic outcomes.

This environment is intentionally lightweight, enabling rapid iteration and reproducibility.

8.3 State Space
Each simulated user is represented by a behavioral state vector:

𝑆
𝑡
=
(
𝐻
𝑡
,
𝑅
𝑡
,
𝜌
)
Where:

𝐻
𝑡
 — habit strength

𝑅
𝑡
 — recency (days since last event)

𝜌
 — responsiveness to incentives

Your uploaded appendix states:

“state definitions, transition functions, reward functions, incentive effects, habit dynamics…”

This chapter formalizes them.

8.4 Behavioral Model
8.4.1 Event Probability
Event probability follows a logistic response:

𝑃
(
𝑒
𝑣
𝑒
𝑛
𝑡
𝑡
=
1
)
=
𝜎
(
𝛼
𝐻
𝑡
+
𝛽
)
Where:

α — sensitivity to habit

β — baseline bias

σ — logistic function

This matches the behavioral model in your simulation code.

8.4.2 Habit Dynamics
Habit evolves according to:

𝐻
𝑡
+
1
=
𝐻
𝑡
+
𝜂
⋅
𝑒
𝑣
𝑒
𝑛
𝑡
𝑡
−
𝛿
Where:

η — habit gain per event

δ — habit decay per day

This is consistent with your Firestore policy fields and simulation code.

8.4.3 Recency Dynamics
𝑅
𝑡
+
1
=
{
0
if 
𝑒
𝑣
𝑒
𝑛
𝑡
𝑡
=
1
𝑅
𝑡
+
1
otherwise
Recency is used for segmentation and fairness analysis.

8.5 Policy Engine
The policy engine mirrors the real Firebase prototype.

8.5.1 Habit Zones
Low habit: 
𝐻
𝑡
<
low threshold

Medium habit: between thresholds

High habit: 
𝐻
𝑡
>
high threshold

8.5.2 Reward Multiplier
𝑅
𝑡
=
𝑏
𝑎
𝑠
𝑒
_
𝑟
𝑒
𝑤
𝑎
𝑟
𝑑
×
𝑓
(
𝐻
𝑡
)
Where:

low habit → boosted multiplier

medium habit → baseline

high habit → soft‑budget reduction

This matches your Firestore policy fields:

soft_budget_factor

responsiveness_weight

habit_low_threshold

habit_high_threshold

8.6 Simulation Loop
For each day 
𝑡
:

Compute event probability

Sample event (0/1)

Compute reward

Update habit

Update recency

Log metrics

This loop is identical to your Python simulation code.

8.7 Experimental Conditions
MIT expects simulation to test multiple conditions:

8.7.1 Baseline Policy (EXP001)
Your default Firestore policy.

8.7.2 Treatment Policies (EXP002, EXP003)
Only added when testing:

higher habit reinforcement

lower reward spend

fairness constraints

budget constraints

8.7.3 Stress Tests
high decay

low responsiveness

high noise

extreme habit thresholds

8.7.4 Sensitivity Analysis
Vary:

η

δ

soft budget

responsiveness

thresholds

8.8 Metrics
Behavioral Metrics
event frequency

return time

habit trajectory

hazard lift

Economic Metrics
total reward spend

cost per incremental event

budget adherence

Fairness Metrics
parity across segments

constraint violations

differential reward allocation

8.9 Validation Criteria
Your uploaded document states:

“Policy outperforms baselines… results are robust across segments… fairness constraints are respected…”

The simulation validates:

effectiveness

efficiency

fairness

robustness

reproducibility

A policy is considered valid if:

it increases event frequency

it maintains or reduces cost per event

it respects budget constraints

it does not disproportionately benefit or penalize segments

8.10 Summary
This chapter defined the full simulation environment for BED v3:

state space

behavioral model

policy engine

simulation loop

experimental conditions

metrics

validation criteria

This simulation is now ready for:

policy comparison

sensitivity analysis

power analysis

integration into the BED v3 Master Document

CTL review
